\section{Bounded uniqueness and connectedness}
\label{sec:bounded-connected-uniqueness}

For a continuous density whose positive set is connected and dense, bounded solutions are unique up to a constant. An energy measure of $w=u-v$ replaces the pointwise gradient calculation available for smooth solutions. Its pushforward under $w$ has a Lipschitz density vanishing at the extreme values. If $w$ were nonconstant, these properties would force both extremal sets to have zero mass for the common Monge--Amp\`ere measure. A smooth test potential with a larger density near one of these sets then contradicts modified comparison.

\begin{theorem}\label{thm:connected-positive-uniqueness}\index{Uniqueness!connected positivity set}
Let $F\in C^0(X)$ satisfy $F\geq0$, $F\not\equiv0$, with $U=\{F>0\}$ connected and dense. If bounded $\omega$-plurisubharmonic functions $u,v$ satisfy
\[
\omega_u^n=\omega_v^n=F\omega^n,
\]
then $u-v$ is constant. In particular, $\sup_Xu=\sup_Xv$ implies $u=v$.
\end{theorem}

Continuity of $F$ supplies the local positive lower bounds and the neighborhoods of the extremal set used below. The auxiliary test densities are smooth even though $F$ is only continuous. Bounded solutions are continuous by \cref{cor:bounded-lp-continuity}, so we may use pointwise extrema throughout the proof. We shall also use modified comparison in the following form. If $b$ is continuous and $\omega_b\geq\varepsilon\omega$ for some $0<\varepsilon<1$, then $b/(1-\varepsilon)$ is $\omega$-plurisubharmonic, so \cref{thm:modified-comparison} applies with $(1-\varepsilon)v=b$. Thus its thin-sublevel estimate applies directly to such a strictly positive comparison potential $b$, and the correction factor tends to one as the height of the sublevel set tends to zero, with $\varepsilon$ fixed.

\begin{proposition}\label{prop:weak-extremum-localization}
Let $u,v$ be continuous $\omega$-plurisubharmonic functions with $\omega_u^n=\omega_v^n=F\omega^n$, where $F\geq0$ is bounded and measurable. Let $q$ be a smooth strictly $\omega$-plurisubharmonic function, and write $\omega_q^n=G\omega^n$. Then for every $0<\kappa<1$ and $\eta>0$,
\begin{equation}\label{eq:weak-localization-min}
F\omega^n\bigl(\{u-v<\min_X(u-v)+\eta\}\cap\{F>\kappa G\}\bigr)>0.
\end{equation}
The analogous assertion holds for the set $\{u-v>\max_X(u-v)-\eta\}$.
\end{proposition}

\begin{proof}
Choose $\delta\in(0,1]$ with $\omega_q\geq\delta \omega$. For $0<t<1$, put $b_t=(1-t)v+tq$, $m_t=\min_X(u-b_t)$, and $E_s=\{u<b_t+m_t+s\}$. Then $\omega_{b_t}\geq t\delta \omega$. Write $\mu=F\omega^n$, $d=\kappa^{-1/n}-1$, and $\gamma=nd$. By \eqref{eq:interior-minkowski-background},
\[
\omega_{b_t}^n\geq\bigl((1-t)F^{1/n}+tG^{1/n}\bigr)^n \omega^n.
\]
\Needspace{9\baselineskip}
On $L=\{F\leq\kappa G\}$ this is at least $(1+\gamma t)\mu$, since $G^{1/n}\geq(1+d)F^{1/n}$ there, and on $R=\{F>\kappa G\}$ it is at least $(1-nt)\mu$. Choose $s>0$ so small that the reverse comparison factor in \eqref{eq:modified-comparison}, with $\varepsilon=t\delta$, is at most $1+\gamma t/2$. Modified comparison on $E_s$, with $x=\mu(E_s\cap L)$ and $y=\mu(E_s\cap R)$, gives
\[
(1+\gamma t)x+(1-nt)y\leq(1+\gamma t/2)(x+y),
\qquad\text{and hence}\qquad y\geq\frac{1-\kappa^{1/n}}2\,\mu(E_s).
\]
The open set $E_s$ is nonempty by the definition of $m_t$. As $\omega_{b_t}^n\geq(t\delta)^n\omega^n$, \eqref{eq:modified-comparison} forces $\mu(E_s)>0$, hence $y>0$.

Let $a=\min_X(u-v)$ and $\Lambda=\|v-q\|_\infty$. Since $u-b_t=u-v+t(v-q)$, we have $|m_t-a|\leq t\Lambda$ and $E_s\subset\{u-v<a+2t\Lambda+s\}$. Given $\eta$, choose $t$ and then $s$ with $2t\Lambda+s<\eta$. Since $y=\mu(E_s\cap R)>0$, this proves \eqref{eq:weak-localization-min}. Interchanging $u$ and $v$ gives the assertion for the near-maximum set.
\end{proof}

\subsection{The energy of the difference}

In this subsection, $u,v$ are continuous $\omega$-plurisubharmonic functions with $A^n=B^n=F\omega^n$ for a bounded measurable function $F\geq0$, and $w=u-v$. Although the cofactor $T$ is now only a current, its product with $\mathrm dw\wedge\dc w$ still defines a positive measure. The identity $\ddc w\wedge T=0$ determines how this energy is distributed over the level sets of $w$.

\Needspace{24\baselineskip}
\begin{proposition}\label{prop:weak-level-energy}
There is a finite positive measure $E=\mathrm{d} w\wedge\dc w\wedge T$, defined below as a limit of regularizations, with the following properties.

For every smooth real function $\chi$ on $\mathbb R$,
\begin{equation}\label{eq:weak-energy-chain}
\int_X\chi''(w)\,\mathrm{d} E=\int_X\chi(w)\,\ddc T.
\end{equation}
The signed current $\ddc T$ is a finite measure. If $a=\min_Xw$, $b=\max_Xw$, and
\[
k(t)=\int_X(w-t)_+\,\ddc T,
\]
then $k$ is nonnegative, globally Lipschitz, supported on $[a,b]$, and
\begin{equation}\label{eq:weak-energy-pushforward}
w_*E=k(t)\,\mathrm{d} t.
\end{equation}
In particular, $E$ charges no level set of $w$.

Let $V$ be a relatively compact coordinate domain such that $F$ has a positive essential lower bound on a neighborhood of $\overline V$. Then there is $C_V<\infty$ such that every piecewise $C^1$ Lipschitz function $p$ on $[a,b]$ satisfies
\begin{equation}\label{eq:weak-local-coercivity}
\int_V|\mathrm{d}(p(w))|\,\mathrm{d}\lambda\leq C_V\left(\int_V|p'(w)|^2\,\mathrm{d} E\right)^{1/2},
\end{equation}
where $\mathrm{d}\lambda$ is Euclidean coordinate volume.
\end{proposition}

\begin{proof}
Locally, expand $T$ into powers of the smooth form $\omega$ multiplied by closed Bedford--Taylor products of $\ddc u$ and $\ddc v$. The operator $\ddc$ then acts only on the powers of $\omega$, and the local mass bounds in \cref{sec:psh-local-capacity} show that $\ddc T$ is a finite signed measure.

We first construct $E$ and prove \eqref{eq:weak-energy-chain}. On a coordinate ball, convolve $u$, $v$ and the coefficients of $\omega$ with a common nonnegative radial mollifier. Write the resulting smooth positive forms as $A_\rho=\omega_\rho+\ddc u_\rho$ and $B_\rho=\omega_\rho+\ddc v_\rho$, and let $T_\rho$ be their mixed sum. On a slightly larger ball, add the same fixed multiple of $|z|^2$ to each potential to obtain plurisubharmonic functions. Their radial regularizations decrease to the originals, so Bedford--Taylor continuity applies to all closed Hessian products. Since $u,v$ are continuous, $u_\rho,v_\rho,w_\rho$ converge uniformly on compact subsets.

For a closed smooth product $S_\rho$ appearing in $T_\rho$, the identity
\begin{equation}\label{eq:weak-energy-regularization}
\mathrm{d} w_\rho\wedge\dc w_\rho\wedge S_\rho
=\tfrac12\ddc(w_\rho^2S_\rho)-w_\rho\ddc w_\rho\wedge S_\rho
\end{equation}
has a distributional limit: both Hessian products on the right converge weakly with uniformly bounded local variation, and their weights converge uniformly. The same holds after multiplication by the convergent smooth factors coming from $\omega_\rho$. Thus the positive measures $E_\rho=\mathrm{d} w_\rho\wedge\dc w_\rho\wedge T_\rho$ converge to a finite positive measure $E$. By \eqref{eq:weak-energy-regularization}, this limit does not depend on the regularization, so the local limits agree on overlapping charts.

For smooth regularizations, the chain rule gives
\[
\ddc\chi(w_\rho)\wedge T_\rho
=\chi''(w_\rho)E_\rho+\chi'(w_\rho)\ddc w_\rho\wedge T_\rho.
\]
The last current tends to zero, because $\ddc w_\rho\wedge T_\rho=A_\rho^n-B_\rho^n$ tends weakly to $A^n-B^n=0$ with uniformly bounded variation, and $\chi'(w_\rho)$ converges uniformly. Passing to the limit and integrating by parts on the compact manifold $X$ gives \eqref{eq:weak-energy-chain}. A partition of unity justifies this passage without a global regularization; the terms containing derivatives of the cutoff functions cancel after summation.

Let $\sigma=w_*(\ddc T)$ and $\nu=w_*E$. Equation \eqref{eq:weak-energy-chain} says that $\nu''=\sigma$ as distributions on $\mathbb R$. Stokes' theorem and $\ddc w\wedge T=0$ give
\[
\sigma(\mathbb R)=\int_X\ddc T=0,\qquad
\int_{\mathbb R}s\,\mathrm{d}\sigma(s)=\int_Xw\,\ddc T=\int_X\ddc w\wedge T=0.
\]
Consequently $k(t)=\int(s-t)_+\,\mathrm{d}\sigma(s)$ is globally Lipschitz and vanishes outside $[a,b]$. It satisfies $k''=\sigma$, so the compactly supported distribution $\nu-k(t)\,\mathrm{d} t$ has second derivative zero and therefore vanishes. Positivity of $\nu$ and continuity of $k$ imply $k\geq0$. This proves \eqref{eq:weak-energy-pushforward}, and the assertion about level sets follows because $\nu$ is absolutely continuous.

We now prove \eqref{eq:weak-local-coercivity}. On a coordinate neighborhood of $\overline V$, set $\beta=\ddc|z|^2$ and choose $c_V>0$ with $F\omega^n\geq c_V\beta^n$. Let $A_{\mathrm{ac}}$ denote the absolutely continuous part of the current $A$, whose coefficients are locally integrable functions. Applying \cref{lem:interior-convolution} with the constant function $g=c_V^{1/n}$ and differentiating the coefficient measures, we obtain $\det_\beta A_{\mathrm{ac}}\geq c_V$ almost everywhere. Bounded plurisubharmonic functions lie in $W^{1,2}_{\mathrm{loc}}$ (\cref{sec:psh-local-capacity}), hence so does $w$ after the quadratic correction above, and its mollifications converge to it strongly in $W^{1,2}_{\mathrm{loc}}$, because the gradient of a mollification is the mollification of the gradient. Thus $w_\rho\to w$ strongly in $W^{1,2}_{\mathrm{loc}}$, and a subsequence of the gradients converges almost everywhere. Moreover, $A_\rho\to A_{\mathrm{ac}}$ almost everywhere, by differentiation of the coefficient measures.

\Needspace{6\baselineskip}
Since $T_\rho$ contains the term $A_\rho^{n-1}$ and all its terms are positive, Fatou's lemma, applied to $E_\rho$, yields the measure inequality
\begin{equation}\label{eq:weak-energy-ac-lower}
E\geq\mathrm{d} w\wedge\dc w\wedge A_{\mathrm{ac}}^{n-1}.
\end{equation}
\Needspace{10\baselineskip}
The cofactor matrix of $A_{\mathrm{ac}}$ has smallest eigenvalue at least $\det A_{\mathrm{ac}}/\operatorname{tr}A_{\mathrm{ac}}\geq c_V/\operatorname{tr}A_{\mathrm{ac}}$, where traces and determinants are taken with respect to $\beta$. The trace is locally integrable because $A$ has locally finite mass. The chain rule for Sobolev functions, \eqref{eq:weak-energy-ac-lower} and the Cauchy--Schwarz inequality give
\[
\int_V|\mathrm{d}(p(w))|\,\mathrm{d}\lambda
\leq C\left(\int_V|p'(w)|^2\,\mathrm{d} E\right)^{1/2}
\left(\int_V\frac{\operatorname{tr}A_{\mathrm{ac}}}{c_V}\,\mathrm{d}\lambda\right)^{1/2}.
\]
The last factor is finite and is absorbed into $C_V$. The values of $p'$ at the finitely many corners of $p$ are immaterial, because $E$ charges no level set of $w$.
\end{proof}

\Needspace{10\baselineskip}
\begin{lemma}\label{lem:weak-null-extremum-levels}
In the setting of \cref{prop:weak-level-energy}, assume in addition that $F$ is continuous and that $U=\{F>0\}$ is connected and dense. If $w=u-v$ is nonconstant, then
\[
F\omega^n(\{w=\min_Xw\})=F\omega^n(\{w=\max_Xw\})=0.
\]
\end{lemma}

\begin{proof}
Put $a=\min_Xw<b=\max_Xw$, and let $k$ be the function of \cref{prop:weak-level-energy}. It is Lipschitz and nonnegative, and $k(a)=0$. Choose $\epsilon_j\downarrow0$ with $a+\epsilon_j<b$. For $\delta>0$, set
\[
J_j(\delta)=\int_a^{a+\epsilon_j}\frac{\mathrm{d} t}{k(t)+\delta}.
\]
If $L$ is a Lipschitz constant for $k$, then $k(t)\leq L(t-a)$ on this interval, so $J_j(\delta)\to\infty$ as $\delta\downarrow0$. Choose $\delta_j>0$ with $J_j=J_j(\delta_j)\geq j$. Let $p_j=0$ on $(-\infty,a]$, $p_j=1$ on $[a+\epsilon_j,\infty)$, and
\[
p_j'(t)=\frac{1}{J_j(k(t)+\delta_j)}\quad(a<t<a+\epsilon_j).
\]
The pushforward identity \eqref{eq:weak-energy-pushforward} gives
\[
\int_X|p_j'(w)|^2\,\mathrm{d} E
=\frac{1}{J_j^2}\int_a^{a+\epsilon_j}\frac{k(t)}{(k(t)+\delta_j)^2}\,\mathrm{d} t
\leq\frac1{J_j}\longrightarrow0.
\]
Every point of $U$ has a coordinate neighborhood $V$ as in \cref{prop:weak-level-energy}, since $F$ is continuous. By \eqref{eq:weak-local-coercivity}, therefore, $\mathrm{d}(p_j(w))\to0$ in $L^1_{\mathrm{loc}}(U)$. The bounded functions $p_j(w)$ converge pointwise and in $L^1_{\mathrm{loc}}$ to $\mathbf1_{\{w>a\}}$, whose distributional gradient therefore vanishes on $U$. Since $U$ is connected, this indicator function is almost everywhere constant on $U$. As $w$ is nonconstant, $\{w>a\}$ is a nonempty open set, and its intersection with the dense open set $U$ is nonempty and open; hence the constant is one. Thus $\{w=a\}\cap U$ has zero volume, and since $F=0$ off $U$, the first assertion follows. Applying the argument to $-w$ gives the assertion for the maximum set.
\end{proof}

\Needspace{18\baselineskip}
\subsection{Concentrating smooth test potentials}

The next proposition produces smooth test potentials whose Monge--Amp\`ere density dominates a fixed multiple of a prescribed function $g$, with a constant independent of $\sup_Xg$ and of the derivatives of $g$. This uniformity allows the test densities to concentrate near a set of measure zero.

\Needspace{10\baselineskip}
\begin{proposition}\label{prop:weak-uniform-test-potentials}
For every $B_0>\|1\|_{L^2(\omega^n)}$ there is $\varepsilon=\varepsilon(X,\omega,B_0)>0$ such that every smooth function $g\geq1$ with $\|g\|_{L^2(\omega^n)}\leq B_0$ admits a smooth strictly $\omega$-plurisubharmonic function $q$ satisfying
\begin{equation}\label{eq:weak-uniform-test-determinant}
\omega_q^n\geq\varepsilon g \omega^n.
\end{equation}
\end{proposition}

\begin{proof}
Choose finitely many coordinate balls $B_j=\{|z_j|<R_j\}$, with closures in their charts, such that the concentric smaller balls $\{|z_j|<r_j\}$ cover $X$. On $B_j$, let $\psi_j$ solve the Dirichlet problem
\begin{equation}\label{eq:weak-local-dirichlet}
(\ddc\psi_j)^n=g\omega^n,\qquad \psi_j|_{\partial B_j}=0.
\end{equation}
The density of $g\omega^n$ with respect to coordinate volume is smooth and strictly positive. The classical smooth Dirichlet theorem therefore applies, since $C_j(|z_j|^2-R_j^2)$ is a smooth strict subsolution with zero boundary values for sufficiently large $C_j>0$. Thus $\psi_j$ is smooth and strictly plurisubharmonic on $\overline B_j$, and the maximum principle gives $\psi_j\leq0$.

We need a lower bound for $\psi_j$ independent of $\sup_Xg$. Write the coordinate equation as $\det(\partial_a\bar\partial_b\psi_j)=D_j$, and put $M_j=-\min_{B_j}\psi_j$. Let $x_j$ be a minimum point. For each real slope $p$ with $|p|<M_j/(2R_j)$, the function $\psi_j-p\cdot(x-x_j)$ attains its minimum at an interior point, because its value at $x_j$ is $-M_j$, while its boundary values are larger than $-M_j$. At such a contact point the real Hessian $S$ of $\psi_j$ is nonnegative, and the gradient of $\psi_j$ maps the contact set onto the ball of these slopes. The area formula and the matrix inequality
\[
\det_{\mathbb R}S\leq4^n\bigl(\det_{\mathbb C}(\partial_a\bar\partial_b\psi_j)\bigr)^2=4^nD_j^2
\]
therefore give $M_j\leq C(n,R_j)\|D_j\|_{L^2(B_j)}^{1/n}$. To prove the matrix inequality, let $J$ be the standard complex structure on the real coordinate space and average $S$ with $J^\mathsf TSJ$. Concavity of $\log\det$, and the fact that the eigenvalues of this $J$-invariant average occur in pairs, give the displayed factor; singular matrices are handled by approximation. Since the ratio of $\omega^n$ to coordinate volume is fixed and bounded, $\|D_j\|_2$ is controlled by $B_0$. Consequently
\begin{equation}\label{eq:weak-local-uniform-sup}
-M_0\leq\psi_j\leq0\quad\text{on every }B_j,
\end{equation}
where $M_0$ depends only on $X,\omega,B_0$ and the chosen finite cover.

Next we glue the local solutions. Put $\beta_j=\ddc|z_j|^2$. Choose $a_j>0$ such that $a_j\beta_j\leq \omega/2$ on $B_j$, and let $t_j=(R_j^2+r_j^2)/2$. Fix $\tau>0$ so small that $\tau M_0<a_j(R_j^2-r_j^2)/4$ for every $j$. The smooth function
\[
s_j=\tau\psi_j+a_j(t_j-|z_j|^2)
\]
is bounded below by a fixed positive constant on the smaller ball, negative on a collar of $\partial B_j$, and satisfies $\omega_{s_j}\geq \omega/2+\tau\ddc\psi_j$. Choose a smooth convex function $\phi_j\colon\mathbb R\to\mathbb R$ approximating $s\mapsto\max(s,0)$, equal to zero below a small negative threshold and to $s$ above the corresponding positive threshold, with $0\leq\phi_j'\leq1$ and $\phi_j''\geq0$; the thresholds are chosen smaller in absolute value than both margins of $s_j$. Set $Q_j=\phi_j(s_j)$ on $B_j$ and extend it by zero. This extension is smooth, and
\[
\omega_{Q_j}=\phi_j'(s_j)\omega_{s_j}+\bigl(1-\phi_j'(s_j)\bigr)\omega
+\phi_j''(s_j)\,\mathrm{d} s_j\wedge\dc s_j\geq \omega/2.
\]
\Needspace{8\baselineskip}
On the smaller ball, $Q_j=s_j$ and $\omega_{Q_j}\geq\tau\ddc\psi_j$. If $N$ is the number of balls, put $q=N^{-1}\sum_jQ_j$. Then $\omega_q=N^{-1}\sum_j\omega_{Q_j}\geq \omega/2$ everywhere. Each point belongs to at least one smaller ball $\{|z_j|<r_j\}$, so positivity of the other summands and \eqref{eq:weak-local-dirichlet} give
\[
\omega_q^n\geq(\tau/N)^n(\ddc\psi_j)^n=(\tau/N)^n g \omega^n.
\]
Thus \eqref{eq:weak-uniform-test-determinant} holds with $\varepsilon=(\tau/N)^n$, independent of $\sup g$ and of all derivatives of $g$.
\end{proof}

\begin{proof}[of \cref{thm:connected-positive-uniqueness}]
By \cref{cor:bounded-lp-continuity}, $u,v$ are continuous. In dimension one, $\ddc w=0$, so the maximum principle makes $w$ constant. Assume $n\geq2$.

Suppose that $w=u-v$ is nonconstant. Let $a=\min_Xw$, $K=\{w=a\}$, and $\mu=F\omega^n$. By \cref{lem:weak-null-extremum-levels}, $\mu(K)=0$. Choose $B_0>\|1\|_{L^2(\omega^n)}$ and let $\varepsilon>0$ be the constant from \cref{prop:weak-uniform-test-potentials}. Since $\mu(K)=0$, the compact set $K_0=K\cap\{F\geq\varepsilon/2\}$ has zero $\omega^n$-volume. Fix $M>1$ with $\varepsilon M>\max_XF$. By outer regularity of the smooth volume measure, choose an open neighborhood $O$ of $K_0$ so small that
\[
\int_X\omega^n+(M^2-1)\int_O\omega^n<B_0^2.
\]
Take a smooth $0\leq\chi\leq1$ equal to one near $K_0$ and supported in $O$, and set $g=1+(M-1)\chi$. If $K_0=\varnothing$, take $\chi=0$. Then $g\geq1$, $g=M$ on $K_0$, and $\|g\|_{L^2(\omega^n)}<B_0$, since $g^2\leq1+(M^2-1)\chi$.

By \cref{prop:weak-uniform-test-potentials}, there is a smooth strictly $\omega$-plurisubharmonic $q$ with $\omega_q^n=G\omega^n\geq\varepsilon g\omega^n$. On $K_0$ we have $G\geq\varepsilon M>F$; on $K\setminus K_0$ we have $F<\varepsilon/2<\varepsilon\leq G$. Thus $G>F$ throughout the compact set $K$. By compactness and continuity, we can choose $0<\kappa<1$ such that $F<\kappa G$ on $K$; put $W=\{F<\kappa G\}$. The complement of this open neighborhood of $K$ is compact, and $w-a$ has a strictly positive minimum there if it is nonempty. Hence some $\eta>0$ satisfies $\{w<a+\eta\}\subset W$.

\Cref{prop:weak-extremum-localization} asserts that $\mu(\{w<a+\eta\}\cap\{F>\kappa G\})>0$. By the preceding inclusion, however, this set is empty. This contradiction shows that $w$ is constant. If $\sup_Xu=\sup_Xv$, the constant is zero.
\end{proof}

\Needspace{16\baselineskip}
For a bounded measurable density, a positive lower bound gives the same conclusion. In this case coercivity holds on all of $X$, and the concentrating test potential can dominate the essential upper bound of the density on a whole neighborhood of the extremal set.

\begin{corollary}\label{cor:bounded-positive-density-uniqueness}
If $0<c_0\leq F\leq C_0<\infty$ almost everywhere and bounded $\omega$-plurisubharmonic functions $u,v$ satisfy $\omega_u^n=\omega_v^n=F\omega^n$, then $u-v$ is constant.
\end{corollary}
\begin{proof}
The functions $u,v$ are continuous by \cref{cor:bounded-lp-continuity}, and in dimension one the assertion follows from $\ddc(u-v)=0$ and the maximum principle. Let $n\geq2$, suppose that $w=u-v$ is nonconstant, and put $K=\{w=\min_Xw\}$. The proof of \cref{lem:weak-null-extremum-levels} applies with $U=X$, since \cref{prop:weak-level-energy} gives the pushforward identity and local coercivity on every chart. Thus $K$ has zero volume.

Fix $B_0>\|1\|_{L^2(\omega^n)}$ and the corresponding $\varepsilon$ in \cref{prop:weak-uniform-test-potentials}. Choose $M>1$ with $\varepsilon M>C_0$, and a smooth cutoff function $\chi$, equal to one on a neighborhood $W$ of $K$ and supported in an open set of volume so small that $g=1+(M-1)\chi$ satisfies $\|g\|_2<B_0$. The resulting test potential $q$ has density $G\geq\varepsilon g$. Choose $\kappa<1$ with $C_0<\kappa\varepsilon M$. Then $F<\kappa G$ almost everywhere on $W$. By continuity of $w$, a sufficiently thin set $\{w<\min_Xw+\eta\}$ lies inside $W$, which contradicts \cref{prop:weak-extremum-localization}. Hence $w$ is constant.
\end{proof}
